% Build: tectonic joint-minimum-problem.tex
\documentclass[10pt]{article}
\usepackage[T1]{fontenc}
\usepackage{amsmath,amssymb}
\usepackage[a4paper,margin=17mm]{geometry}
\usepackage{microtype}
\usepackage[hidelinks]{hyperref}
\pagestyle{empty}
\setlength{\parindent}{0pt}
\setlength{\parskip}{5pt}
\newcommand{\Z}{\mathbb Z}
\newcommand{\Q}{\mathbb Q}
\newcommand{\B}{\mathcal B}
\newcommand{\E}{\mathcal E}
\newcommand{\F}{\mathcal F}
\begin{document}
\begin{center}
{\Large\bfseries No-three-collinear lattice walks: the joint minimum problem}\\[3pt]
{\small A sharp-threshold follow-up to Erd\H{o}s Problem 193}
\end{center}

\textbf{Definitions.}
A finite or infinite lattice walk $(X_n)$ is \emph{admissible} if $X_0=0$,
its vertices are pairwise distinct, and
\[
 \operatorname{rank}_{\Q}[\,X_j-X_i\quad X_k-X_j\,]=2
 \qquad\text{for every }0\le i<j<k
\]
within its index range. Thus no three vertices lie on an affine line.
For $d,s\ge1$, define
\begin{align*}
 \B(d)&\iff \exists\text{ an infinite admissible }(P_n)\subset\Z^d:
       \quad P_{n+1}-P_n\in\{e_1,\ldots,e_d\}\quad(n\ge0),\\
 \E(3,s)&\iff \exists S\subset\Z^3\setminus\{0\},\ |S|\le s,\
       \exists\text{ an infinite admissible }(Q_n)\subset\Z^3:\quad
       Q_{n+1}-Q_n\in S\quad(n\ge0),\\[-2pt]
 d_*&=\min\{d\ge1:\B(d)\},\qquad
 s_*=\min\{s\ge1:\E(3,s)\}.
\end{align*}
Here $e_r$ is the $r$th positive standard basis vector. In $\E(3,s)$,
$S$ is chosen \emph{once for the entire infinite walk}; its vectors are
arbitrary nonzero integer vectors, with no prescribed coordinate bound.
Both minima are finite by the original finite-step construction and basis encoding.

\textbf{Problem.}
Determine the exact ordered pair $(d_*,s_*)$: prove unconditional infinite
existence at each claimed minimum and impossibility for \emph{every} walk
with a smaller dimension or step budget, respectively.
In particular, determine whether the compression gap $s_*-d_*$ is zero.
Equality alone would not determine the common value.

\textbf{Basic reduction and bounds.}
Label the vectors of a successful step set by $1,\ldots,t$, $t\le s$,
and replace each step of type $r$ by $e_r$. The resulting count walk
projects linearly onto the original walk, so is itself admissible. Hence
\[
 \E(3,s)\Longrightarrow\B(s),\qquad
 \boxed{4\le d_*\le s_*\le16}.
\]
The lower bound follows because every ternary word of length eight contains
an abelian square; the ceiling comes from the original sixteen-step 3D
construction. The converse implication $\B(s)\Rightarrow\E(3,s)$ is
\emph{not assumed}. Equivalently, $d_*$ is the minimum number of step types
when the ambient lattice dimension is unrestricted.
A written six-step construction proposes $s_*\le6$, hence $d_*\le6$,
but independent review remains pending. Subject to that verification, the
candidate pairs are $(4,4),(4,5),(4,6),(5,5),(5,6),(6,6)$; none is asserted
as the answer.

\textbf{Equivalent word and projection formulations.}
Let $\mathcal W_d$ consist of infinite words on $\{1,\ldots,d\}$ having no
adjacent nonempty blocks $x,y$ with
$\psi(x)/|x|=\psi(y)/|y|$, where $\psi$ is the Parikh (letter-count) vector.
The lengths $|x|,|y|$ may differ: avoiding ordinary abelian squares alone
is insufficient. If $P_n(w)$ counts the first $n$ letters, then
\[
 \B(d)\iff\mathcal W_d\ne\varnothing,\qquad
 \E(3,d)\iff
 \exists w\in\mathcal W_d\ \exists A\in\Z^{3\times d}\ \forall i<j<k:
 \operatorname{rank}_{\Q}
 [\,A(P_j-P_i)\quad A(P_k-P_j)\,]=2.
\]
Thus $s_*=d_*$ requires a fixed integer projection for \emph{one}
optimal-alphabet witness, not for every valid word.
For basis walks only, $\sum_r(P_n)_r=n$ makes noncollinearity equivalent to
$(k-j)(P_j-P_i)\ne(j-i)(P_k-P_j)$.
This time-weighted test is not sufficient for arbitrary 3D step sets.

\textbf{Equivalent quantifier distinction.}
Let $\F_N(S)$ mean that an admissible $N$-step walk in $\Z^3$ uses steps
from $S$. In the following formulas, menus are finite subsets of
$\Z^3\setminus\{0\}$:
\[
 \boxed{\begin{aligned}
 \B(s)&\iff \forall N\ge1\ \exists S_N,\ |S_N|\le s:\ \F_N(S_N),\\
 \E(3,s)&\iff \exists S,\ |S|\le s\ \forall N\ge1:\ \F_N(S).
 \end{aligned}}
\]
These follow from integer projection of finite basis prefixes, basis
encoding, and K\"onig's lemma. A uniform bound on all menu coordinates in
the first line suffices to obtain the second; unbounded prefix-dependent
projections do not. Finite positive evidence and lower bounds confined to
one construction family do not settle either minimum.

\vfill
{\footnotesize\textit{Provenance and status (6 September 2026).}
Original finite-step theorem: Stijn Cambie and Erik Kalviainen;
positive-basis/word formulation and encoding: Jeffrey Shallit;
proposed six-step improvement: Kalviainen, using Cambie's offsets and
Shallit's encoding. This is a research problem statement, not a new joint
authorship claim or a resolution of these follow-up minima.}
\end{document}
